\section*{Chapter \ref{ch:b2:plant-plasticity} --- Plant Adaptations and Phenotypic Plasticity}

\begin{solution}{exo:b2:plant-plasticity:1}
Plasticity: one \omterm{def:b2:meiosis-heredity:vocabulary}{genotype}, several \omterm{def:b2:meiosis-heredity:vocabulary}{phenotypes} according to the
environment (sun and shade leaves on one oak). \omterm{def:b2:plant-plasticity:plasticity}{Norm of reaction}:
the curve of a \omterm{def:b2:meiosis-heredity:vocabulary}{genotype}'s \omterm{def:b2:meiosis-heredity:vocabulary}{phenotype} against an environmental
variable (leaf thickness against light). \omterm{def:b2:plant-plasticity:plasticity}{Acclimation}: a reversible
physiological adjustment (frost hardening of rye in autumn).
\omterm{def:b2:plant-plasticity:plasticity}{Adaptation}: a heritable fit selected over generations (the CAM
metabolism of a cactus).
\end{solution}

\begin{solution}{exo:b2:plant-plasticity:2}
Intact coleoptile lit from one side: bends toward the light. Tip
covered with an opaque cap: no bending --- the tip perceives the
light. Tip covered with a transparent cap: bends --- the cap itself
is harmless. Base shielded by an opaque tube, tip exposed: bends ---
the perception is at the tip and the response below it.
\end{solution}

\begin{solution}{exo:b2:plant-plasticity:3}
\omterm{prop:b2:plant-plasticity:sunshade}{Sun leaf}: small, thick, two or three palisade layers, thick
cuticle, many stomata, high maximal photosynthesis and high
compensation point. \omterm{prop:b2:plant-plasticity:sunshade}{Shade leaf}: large, thin, one palisade layer,
more chlorophyll per gram, low respiration, low compensation point.
The choice is made while the leaf is a primordium in the bud, from
the light it receives; a mature leaf cannot change.
\end{solution}

\begin{solution}{exo:b2:plant-plasticity:4}
The \omterm{prop:b2:plant-plasticity:stomata}{guard cells} take up potassium and water, swell, and, their
inner walls being thicker, bow apart to open the pore. Blue light
and low internal $\mathrm{CO_2}$ open them in the morning; \omterm{def:b2:plant-flowering:dormancy}{abscisic
acid}, from roots in drying soil or from the leaf, closes them in a
drought, as does very dry air.
\end{solution}

\begin{solution}{exo:b2:plant-plasticity:5}
$\varphi = R/(R + 0.77\,FR)$ with $FR = 1$: $1.15$: $0.60$; $0.7$:
$0.48$; $0.2$: $0.21$; $0.05$: $0.06$. The threshold $0.4$ falls
between the gap ($0.48$) and the canopy ($0.21$): \omterm{thm:b2:plant-plasticity:phytochrome}{shade avoidance}
switches on under the canopy.
\end{solution}

\begin{solution}{exo:b2:plant-plasticity:6}
$E = 0.25\times 0.025 = \qty{6.25}{mmol.m^{-2}.s^{-1}}$; $A =
0.25\times 100/1.6 = \qty{15.6}{\micro mol.m^{-2}.s^{-1}}$; $A/E =
2.5\times 10^{-3}$ mol of carbon per mol of water, i.e.\ 400 water
molecules per carbon, $400\times 18/12 = \qty{600}{g}$ of water per
gram of carbon. In 12 hours: $6.25\times 10^{-3}\times 43200 =
\qty{270}{mol/m^2}$, \qty{4.9}{L} per square metre.
\end{solution}

\begin{solution}{exo:b2:plant-plasticity:7}
$E = 0.25\times 0.005 = \qty{1.25}{mmol.m^{-2}.s^{-1}}$, $A$
unchanged: $A/E = 0.0125$, 80 waters per carbon, \qty{120}{g/g} ---
a fifth of the daytime cost.
\end{solution}

\begin{solution}{exo:b2:plant-plasticity:8}
$v = 2r^{2}\Delta\rho g/9\eta = 2\times 4\times 10^{-12}\times
400\times 9.81/0.09 = \qty{3.5e-7}{m/s}$: \qty{0.35}{\micro m/s},
about \qty{30}{s} to cross the cell. On a clinostat the direction of
gravity relative to the cell changes every few seconds of settling,
so the grains never accumulate on one side and the averaged
stimulus over the plant's integration time (minutes) is zero.
\end{solution}

\begin{solution}{exo:b2:plant-plasticity:9}
Mean extension in two hours $0.04$; sides $1.2\times 0.04 = 0.048$
and $0.8\times 0.04 = 0.032$; $\theta = 20\times 0.016/2 = 0.16$
rad, about \qty{9}{\degree}. Bending accumulates at $0.08$ rad per
hour; $\pi/2$ takes about \qty{20}{h}.
\end{solution}

\begin{solution}{exo:b2:plant-plasticity:10}
Reserves allow \qty{50}{cm} of stem. Nettle bed: \qty{30}{cm} in
7.5 days for \qty{30}{mg} --- the seedling reaches light with
reserves to spare. \omterm{def:b2:plant-meristems:wood}{Wood}: \qty{10}{m} is out of reach; the seedling
spends its \qty{50}{mg} on half a metre of pale stem in twelve
days and dies etiolated, whereas a shade-tolerant strategy --- a
few thin leaves, slow growth --- might have kept it alive for years.
\end{solution}

\begin{solution}{exo:b2:plant-plasticity:11}
Ice forms first in the extracellular spaces, where the solution is
more dilute; the vapour pressure over ice is lower than over the
cell's water, so water leaves the cell to join the ice and the cell
dehydrates and its membranes collapse. A week of cold loads the
cell with sugars and protective proteins that lower its freezing
point and stabilise membranes, and makes antifreeze proteins that
limit crystal growth. Keeping the programme on all year diverts
carbon from growth to protection and slows division: a permanently
hardened plant is a small one.
\end{solution}

\begin{solution}{exo:b2:plant-plasticity:12}
The far-red signal is honest when the neighbour can be overtopped
and a lie when it is a tree; the warm week is honest in April and
a lie in February; in a dense crop every plant's neighbours are its
equals, so the signal is honest in form but useless, and breeders
select plants that ignore it. A \omterm{def:b2:plant-plasticity:plasticity}{norm of reaction} is the record of
how often each cue has told the truth in the lineage's past, and it
fails whenever the present differs from that past.
\end{solution}

\begin{solution}{pb:b2:plant-plasticity:1}
\textbf{1.} Open: $1.15/(1.15 + 0.77) = 0.60$; canopy: $0.2/(0.2 +
0.77) = 0.21$.
\textbf{2.} Low $\varphi$: plants overhead --- elongate. Under a
shade cloth $\varphi$ stays $0.60$: a dark day, not a neighbour; no
\omterm{thm:b2:plant-plasticity:phytochrome}{shade avoidance}, only slower growth for want of light.
\textbf{3.} $1 + 2(0.60 - 0.21) = 1.78$: \qty{2.7}{cm} a day.
\textbf{4.} \qty{37}{cm} against \qty{21}{cm}.
\textbf{5.} Bending toward the gap (phototropism, via the blue-light
receptor phototropin) and faster elongation and leaf reorientation
on the side with the higher $R{:}FR$ (via phytochrome).
\textbf{6.} Shade leaves are built large and thin to catch scarce
light cheaply; in full sun they overheat, photo-inhibit and scorch,
and the plant must replace them with sun leaves.
\textbf{7.} \qty{2.7}{cm} a day over a \qty{15}{mm} zone is
\qty{0.11}{cm} an hour, a relative extension of $0.074$ per hour;
shaded side $1.3\times 0.074 = 0.096$, lit side $0.7\times 0.074 =
0.052$.
\textbf{8.} $\theta = 15\times(0.096 - 0.052)/2 = 0.33$ rad, about
\qty{19}{\degree} in an hour.
\textbf{9.} \qty{60}{\degree} is $1.05$ rad: about three hours.
\textbf{10.} $v = 2\times 6.25\times 10^{-12}\times 400\times
9.81/0.09 = \qty{5.5e-7}{m/s}$, \qty{0.55}{\micro m/s}: \qty{22}{s}
to cross the cell.
\textbf{11.} Phototropism, and \omterm{thm:b2:plant-plasticity:phytochrome}{shade avoidance} itself raises the
angle at which the shoot is content to grow: light is the resource
being fought for, and a stem that straightened against gravity would
turn back into the shade.
\textbf{12.} The root bends downward until it is vertical: statoliths
settle to the lower side, auxin accumulates there, and in a root the
extra auxin inhibits growth, so the upper side grows faster. Same
redistribution, opposite sensitivity of the cells.
\textbf{13.} $E = 0.2\times 0.015 = \qty{3}{mmol.m^{-2}.s^{-1}}$; over
$2\times 10^{-3}$ \unit{m^2} and \qty{50400}{s}: \qty{0.30}{mol},
\qty{5.4}{g} of water a day.
\textbf{14.} $A = 0.2\times 80/1.6 = \qty{10}{\micro mol.m^{-2}.s^{-1}}$;
per day $0.50$ mol of carbon per square metre; $A/E = 3.3\times
10^{-3}$: 300 waters per carbon, \qty{450}{g/g}.
\textbf{15.} $0.50\times 2\times 10^{-3} = \qty{1}{mmol}$: \qty{12}{mg}
of carbon a day.
\textbf{16.} Gap: $E = \qty{6}{mmol.m^{-2}.s^{-1}}$, \qty{10.9}{g} a
day; efficiency \qty{900}{g/g}.
\textbf{17.} $E = 0.02\times 0.03 = \qty{0.6}{mmol.m^{-2}.s^{-1}}$
(tenfold down); $A = 0.02\times 200/1.6 = \qty{2.5}{\micro
mol.m^{-2}.s^{-1}}$ (fourfold down). Water falls more: as the leaf
draws down its internal $\mathrm{CO_2}$ the gradient for carbon
widens while that for water cannot.
\textbf{18.} \qty{4.25}{g} of water, \qty{0.85}{g} to lose; at
\qty{10.9}{g} a day ($0.78$ g an hour) it wilts in about an hour.
\textbf{19.} CAM is slow and needs succulent storage tissue; a
seedling racing for light under a canopy, where water loss is small
anyway, needs speed, not economy.
\textbf{20.} $100/2 = \qty{50}{cm}$.
\textbf{21.} Nettle bed: \qty{40}{cm} in 15 days for \qty{80}{mg} ---
it reaches light. \omterm{def:b2:plant-meristems:wood}{Wood}: \qty{8}{m} is unreachable; it spends its
reserves on half a metre and dies. Better in the wood: stay short,
make thin shade leaves and wait for a gap --- the beech's way.
\textbf{22.} Birch: a steep norm, elongating strongly at low
$\varphi$, since a pioneer's neighbours are its own size. Beech: a
flat norm, little elongation and much tolerance, since its neighbours
are trees.
\textbf{23.} Disable the shade-avoidance response --- keep
phytochrome B active or remove the \omterm{prop:b2:cell-differentiation:transcriptional}{transcription factors} it
restrains --- and check flowering time, germination and stem
strength, which the same pathway controls.
\textbf{24.} A plant that cannot sense neighbours is overtopped
whenever it could have won; a plant that always believes the cue
squanders itself under every tree. The winning norm bets on the
cue in proportion to how often, in the lineage's past, it was
right.
\textbf{25.} $\varphi = 0.60$ in the open and $0.21$ under the canopy;
\qty{37}{cm} after 14 days; \qty{5.4}{g} of water a day in the shade,
\qty{10.9}{g} in the gap.
\end{solution}
